\documentclass[a4paper,12pt]{article}
\usepackage[utf8]{inputenc}
\usepackage{graphicx}
\usepackage{fancyhdr}
\usepackage{tabularx}
\usepackage{ragged2e}
\usepackage[table]{xcolor}
\usepackage{float}

\renewcommand{\abstractname}{}
\renewcommand{\arraystretch}{1.5}
\newcolumntype{Y}{>{\RaggedRight\arraybackslash}X} 

\DeclareGraphicsExtensions{.png}
\begin{document}
\pagestyle{fancy}
\fancyhead{}
\lhead{\includegraphics[height=20px]{epic}}
\rhead{Project Plan}

\setcounter{secnumdepth}{0}

\author{Team BananaCore\\Product Development Team 35/36}
\title{Project Plan}
\date{February 2016}

\maketitle
\begin{abstract}
	\begin{tabular}[h]{rl}
		Project's full name: & Economic Performance Information Console\\
		Project's short name: & EPIC \\
		Last modified: & 03.02.2016
	\end{tabular}
	\newline
	\vspace{20px}
	\newline
	\begin{tabularx}{\textwidth}{rllY}
		\hline
		\multicolumn{4}{c}{Change History: \cellcolor[gray]{0.9}} \\
		 \cellcolor[gray]{0.9} Version & \cellcolor[gray]{0.9} Date \cellcolor[gray]{0.9} & \cellcolor[gray]{0.9} Made By & \cellcolor[gray]{0.9} Comments  \\
		\hline
		0.1 & 03.02.2016 & Ole-Magnus Pedersen & First draft, for first hand-in. \\
		\hline
		0.1.1 & 04.02.2016 & Eivind Grimstad & Added time spent and DoD. \\
		\hline
		0.2 & 19.02.2016 & Eivind Grimstad & Added more roles. \\
		\hline
		0.3 & 04.03.2016 & Ole-Magnus Pedersen & Updated after sprint 1 \\
		\hline
		0.4 & 18.03.2016 & Ole-Magnus Pedersen & Updated after sprint 2 \\
		\hline
		0.5 & 08.04.2016 & Ole-Magnus Pedersen & Updated after sprint 3 \\
		\hline
	\end{tabularx}
\end{abstract}

\newpage
\tableofcontents
\newpage

\section{Introduction}
This document presents the framework and methods used in the development of this product, EPIC. It was drawn up by the Product Development Team together, spending about 40 man-hours.


\section{Development method}
We will be using an agile method where eXtreme programming/Scrum is the basis. Our backlog will use KanBan, scheduling is inspired by Mobile-D.

\section{Aims}
EPIC is a system that aims to inform the driver about his or her driving performance. It measures the efficiency of the driver’s braking and presents data both live and in graph form. In addition it conveys information about fuel usage, engine rpm and speed, and recommends gear shifting for efficient driving.

\section{Results}
For each sprint, the following should be produced:
\begin{itemize}
	\item A working system with each release
	\item Executables (Java Source Code)
	\item Documentation of the last sprint
\end{itemize}

\section{Tasks}
As we are using an agile method of development (see introduction), we will be using the backlog to keep track of tasks as well as their date of completion.

\subsection{Sprint 1}
\begin{figure}[H]
	\centering
	\includegraphics[width=\textwidth]{sprint1_backlog}
	\caption{Backlog for sprint 1}
\end{figure}

\begin{figure}[H]
	\centering
	\includegraphics[width=\textwidth]{sprint1_burndown}
	\caption{Burndown chart for sprint 1}
\end{figure}

\subsection{Sprint 2}
\begin{figure}[H]
	\centering
	\includegraphics[width=\textwidth]{sprint2_backlog}
	\caption{Backlog for sprint 2}
\end{figure}

\begin{figure}[H]
	\centering
	\includegraphics[width=\textwidth]{sprint2_burndown}
	\caption{Burndown chart for sprint 2}
\end{figure}

\subsection{Sprint 3}
\begin{figure}[H]
	\centering
	\includegraphics[width=\textwidth]{sprint3_backlog}
	\caption{Backlog for sprint 3}
\end{figure}

\begin{figure}[H]
	\centering
	\includegraphics[width=\textwidth]{sprint3_burndown}
	\caption{Burndown chart for sprint 3}
\end{figure}

\section{Definition of Done}
A user story is defined as done after the features in it have been implemented,
the code has been properly documented where necessary and it has been
unit tested. Integration testing will be done on a higher level (between each
sprint), and is not needed for a story to be defined as done.

\section{Schedule}
\begin{tabularx}{\textwidth}{YYY}
	\hline
	Weeks -- deadline & Phase & Deliverables \\ \hline 
	3-5 -- 5th February & Exploration & Concept, Project Plan, Backlog \\
	6-7 -- 19th February & Initialization & Set up Git, Google Drive, Raspberry Pi and 3D-print a RPI-case. \\
	8-9 -- 18th March & Sprint 1 & Implement skeleton/basics of system \\
	10-11 -- 23rd March & Sprint 2 & Keep going down backlog, implement more functionality from backlog. \\
	12 -- 21st to 23rd March & User testing & Somewhat functional MVP for testing \\
	13-14 -- 8th April & Sprint 3 & Last functionality of MVP \\
	15-16 -- 11th to 22nd April & Stabilization & Stable MVP \\
	17 -- 25th to 28th April & Delivery & MVP \\ \hline
\end{tabularx}

\section{Estimated cost}
We estimate 8 persons working 8 hours each for 12 weeks. 768 total hours estimated.

In addition to this, there is some hardware cost (price of the Pi and Screen).

\section{Completion of the project}
Resources (hours) are fixed, with teams consisting of 2 * 4 students each spending 8 hours per week. This adds up to a total of 12 weeks * 8 student * 8 hours = 768 hours of work. Whatever system is finished by this date, will be considered the finished product. Hopefully a working MVP.

\section{Risk Assesment}
\begin{tabularx}{\textwidth}{YYY}
	\hline
	Risk & Means to Prevent & Responsibility \\
	\hline
	Meeting absence & Have to bring cake to the next meeting if you are absent & Everyone\\
	Major bugs & Early and frequent testing & Everyone \\
	Conflicting schedules & Plan meetings early. Communicate & Product owner for Dev Team meetings and Feature Lead for Feature Team meetings. \\
	Technical problems & Peer-programming and unit testing & Everyone \\ \hline
\end{tabularx}

\section{Quality Assurance}
Following an agile development method, we intend to focus on the acceptance criteria in the backlog as quality. We want an MVP, not a perfect product as our time is short. However, we will be using unit tests (JUnit with Java). 
There is also a planned user-test in march, where we will simulate data to see if we have a working product.

\section{Documentation}
\begin{itemize}
	\item Monitoring of time, on a daily basis
	\item Meeting summaries, both as full product development teams and feature teams
	\item Documentation of system, updated after each sprint
\end{itemize}

\section{Resources and Organization}
\begin{tabularx}{\textwidth}{p{.4\textwidth}p{.6\textwidth}}
	Project Development Team Leader and Product Owner & Eivind Grimstad \\
	Feature Team 36 - Quad Core & 
		Eivind Grimstad (leader)\newline
		Ole-Magnus Pedersen (GUI responsible)\newline
		Marton Skjæveland (scrum master)\newline
		Andreas Stensbye (lead architect)\newline
		\\
	Feature Team 35 - Banana & 
		Sander Aker Christiansen (leader, lead tester)\newline
		Carlo Alfredo Morte III\newline
		Henrik Bjelke Anderson (scrum master)\newline
		Ida Rostveit
\end{tabularx}
\newline
\vspace{10px}

The Product Development Team of 8 is divided into two Feature Teams of four people each. During each sprint these teams will work independently on a designated part of the system.

After the sprint the Product Development Team will come together to integrate the features they have implemented.

\section{Technical Requirements}
\begin{itemize}
	\item Software stack:
	\begin{itemize}
		\item Java 8 with JavaFX for GUI
		\item JUnit for testing
		\item MySQL for database management
	\end{itemize}
	\item Hardware:
	\begin{itemize}
		\item Raspberry Pi Model 2
		\item 3.5" touch-sensitive monitor
	\end{itemize}
	\item Version Control
	\begin{itemize}
		\item GIT
	\end{itemize}
	\item IDE
	\begin{itemize}
		\item IntelliJ IDEA
	\end{itemize}
\end{itemize}

\section{Appendix A: Retrospective analysis}
\subsection{Sprint 1}
Positive feedback:
\begin{itemize}
	\item Effective when working, well defined boundary of working time
	\item Good sprint planning, finished with sprint backlog
	\item Good to have the same room every time
\end{itemize}
 Negative feedback:
 \begin{itemize}
 	\item Some downtime while integrating
 	\item Backlog items were dependent on each other. Made it hard to test
 \end{itemize}
 Change:
 \begin{itemize}
 	\item New partners? Agreed it is not necessary, but could be done.
 	\item It will probably be easier to work independently now that we have a product skeleton
 \end{itemize}
 
 \subsection{Sprint 2}
 Positive feedback:
 \begin{itemize}
 	\item When we meet up, we do pretty effective work
 \end{itemize}
 Negative feedback:
 \begin{itemize}
 	\item We overestimate how much we are able to do.
 \end{itemize}
 Change:
 \begin{itemize}
 	\item Be more clear about who does what during sprint planning, the settings were a confusing part of this sprint.
 	\item Be more realistic about how much we are able to do (as we are nearing the point where we NEED an actual functioning MVP).
 	\item Better planing for the retrospective meeting\
 \end{itemize}
 
 \subsection{Sprint 3}
 Positive feedback:
 \begin{itemize}
 	\item Good communication between the different pairs
 \end{itemize}
 Negative feedback:
 \begin{itemize}
 	\item Worked to little to finish all tasks
 	\item Didn't update kanban board when starting/finishing tasks
 \end{itemize}
 Change:
 \begin{itemize}
 	\item The pairs should plan better when they work
 	\item Use more time on administrative task (i.e. retrospective meetings, finishing deliverables, planning user tests)
 	\begin{itemize}
 		\item Retrospective meetings at 13:00 the last wednesday of the sprint
 	\end{itemize}
 \end{itemize}
 Start:
 \begin{itemize}
 	\item Write unit tests
 \end{itemize}
\end{document}